\documentclass[11pt]{article}
\usepackage[margin=1in]{geometry}
\usepackage{amsmath,amssymb}
\usepackage[T1]{fontenc}
\usepackage{lmodern}

\title{Literature-implied answer: a type-(2) HI sequentially minimal space}
\author{arXiv:1104.4724 answered in scope by arXiv:1112.2411}
\date{2026-07-03}

\begin{document}
\maketitle

\section*{Status}

\textbf{literature\_implied\_answer (final Problem (a) / type-(2) example only).}

\section*{Original Question}

The source paper is V. Ferenczi and C. Rosendal, \emph{Banach spaces without
minimal subspaces - Examples}, arXiv:1104.4724.  In the final section
``Problems and comments'', the authors ask for examples filling missing
classes in their classification chart.  One listed subproblem is:
\[
  \text{Find a HI space which is sequentially minimal.}
\]
Equivalently, this asks for an example in the missing type-(2) row:
\[
  \text{HI, tight, sequentially minimal.}
\]

\section*{Supporting Literature}

The supporting paper is V. Ferenczi and Th. Schlumprecht,
\emph{Subsequential minimality in Gowers and Maurey spaces},
arXiv:1112.2411.  Its abstract states that the authors define block sequences
in every block subspace of a variant of the Gowers-Maurey space so that the
map \(x_{2n-1}\mapsto x_{2n}\) extends to an isomorphism, and that this
implies the existence of a subsequentially minimal HI space.

In the introduction, the supporting paper states the result explicitly:
\[
  \text{There exists a tight, HI, sequentially minimal space } \mathcal X_{GM}.
\]
The final construction is stated as Theorem 6.15 and Theorem 6.16 in the
source TeX, where a tight HI block subspace \(\mathcal X_{GM}\) of a
Gowers-Maurey variant is obtained with a normalized basis that is
subsequentially minimal.

\section*{Conclusion and Scope}

Therefore arXiv:1112.2411 supplies the requested HI sequentially minimal
example and fills the type-(2) entry left open in arXiv:1104.4724.

This does not settle the other final problems in arXiv:1104.4724: it does not
produce type (5b) or type (5d) examples, does not answer the tight-with-
constants uniformly inhomogeneous questions, and does not determine whether
the Argyros-Haydon space is of type (1) or type (2).

\section*{Search Evidence}

The run indexes contained no exact prior packet for arXiv:1104.4724.  Searches
for the core phrases ``HI sequentially minimal'', ``type (2)'', and
``Gowers-Maurey'' identify arXiv:1112.2411 as the decisive source.  Later
papers arXiv:1206.0651 and arXiv:1303.2370 also record that the
Ferenczi-Schlumprecht construction completes the HI tight sequentially minimal
class.

\end{document}
